% =====================================================================
% Exemple complet — TD — thème `ocots`
%
% Les énoncés se collent tels quels dans le polycopié : c'est le point de la
% refonte. Le corps est dans ../../content/td-body.tex, partagé par les 4
% thèmes.
% =====================================================================
\documentclass[11pt]{ocots-td}
\usepackage[lang=fr, theme=ocots, solutions=none, math=analysis,
            institution={n7}]{ocots}
\lstset{language=[LaTeX]TeX}

\title{Équations différentielles linéaires}
\shorttitle{TD 1 --- Équations linéaires}
\numero{TD n\textsuperscript{o}~1}
\date{2025--2026}
\discipline{Calcul différentiel et EDO}
\promotion{ENSEEIHT / INSA --- 1A}

\begin{document}

\maketitle

\begin{lstlisting}
\usepackage[..., theme=ocots, ...]{ocots}
\end{lstlisting}
\begin{quote}\itshape
Ce TD illustre le th\`eme \texttt{ocots} : m\^eme \'enonc\'e que pour les
autres th\`emes (\texttt{examples/themes/}) ; seule cette ligne change.
\end{quote}

\input{../../content/td-body}

\end{document}
